% ===========================================================================
%  Methodology section for the INFOCOM systems paper.
%  Requires: \usepackage{booktabs}, \usepackage{amsmath}, \usepackage{siunitx}
%  (or replace \si/\num macros with plain text if siunitx is unavailable).
% ===========================================================================

\section{System Design and Methodology}
\label{sec:methodology}

We describe a GPU-resident, real-time wideband signal-detection system that
ingests \SI{500}{\mega\hertz} of instantaneous bandwidth from a software-defined
radio and produces per-frame time--frequency detection masks. We first present
the end-to-end streaming architecture (\S\ref{sec:arch}), then the coherent-power
detector as a \emph{hybrid image/power} algorithm (\S\ref{sec:detector}), and
finally the labeled evaluation dataset (\S\ref{sec:dataset}) and offline scoring
protocol (\S\ref{sec:eval}).

% ---------------------------------------------------------------------------
\subsection{End-to-End Streaming Architecture}
\label{sec:arch}

The pipeline is implemented as an NVIDIA Holoscan operator graph (single-channel
coherent-power configuration with calibrated per-frequency thresholding). Each edge
carries a device-resident tensor together with the CUDA stream on which it was
produced, so the entire hot path---from network interface to detection mask---
executes on the GPU with no host staging of raw samples. The active operator
chain for the single-channel coherent-power configuration is:
%
\[
\text{CHDR}\!\rightarrow\!\text{FFT}\!\rightarrow\!\text{Detector}
\;\rightarrow\;\{\text{Mask preview},\ \text{Spectrogram preview},\ \text{Logger}\}.
\]

\paragraph{(1) RF ingest and GPUDirect DMA.}
IQ samples are streamed as UHD CHDR/VITA packets over UDP and captured by a
DPDK-managed network stack with flow isolation. The sample payload is DMA'd
\emph{directly into GPU device memory} (a \texttt{device}-kind DPDK memory
region), i.e.\ GPUDirect zero-copy from the NIC; only the small (\SI{42}{}/
\SI{64}{}-byte) packet headers are placed in CPU memory. A CHDR-converter
operator performs on-GPU endianness conversion and casts interleaved
\num[]{16}-bit integer IQ to single-precision complex, normalizing by
$1/\text{0x7FFF}$, and assembles packets into a per-batch device tensor of shape
$[\,N_{\text{ffts}}\!=\!512,\ L_{\text{fft}}\!=\!20480\,]$. Four concurrent CUDA
streams pipeline batch DMA. The operator emits a $512\times20480$ complex tensor
per frame and attaches metadata (channel index, emit timestamps, sample rate,
center frequency, and a partial-batch / soft-resync epoch used to reset detector
state after packet loss).

\paragraph{(2) Wideband FFT.}
The FFT operator runs \num{512} batched length-\num{20480} complex-to-complex
transforms via cuFFT on the inbound stream, then re-centers DC per row
(\texttt{fftshift}). Crucially, the FFT emits \emph{complex} bins---no magnitude
or log conversion happens here; dB formation is deferred to each consumer. The
output is a $512\times20480$ complex time--frequency frame in which rows index
time (\num{512} FFT frames) and columns index frequency (\num{20480} bins at
\SI{24414}{\hertz}/bin spanning \SI{500}{\mega\hertz}). Frame-number and
frequency-axis metadata (\texttt{span}, \texttt{resolution}, bin indices) are
propagated.

\paragraph{(3) Detector.}
The coherent-power detector consumes the complex FFT frame directly (there is no
intermediate spectrogram operator in this configuration) and emits a
device-resident binary detection mask (\S\ref{sec:detector}). Its output port is
unconditioned, so detection never blocks upstream ingest. Partial CHDR batches
are dropped and detector state is reset on a soft-resync epoch to prevent
stale-history contamination after loss.

\paragraph{(4) Visualization and logging.}
A preview operator forms the \emph{display} spectrogram---this is where the
grayscale dB image is computed ($10\log_{10}(\lvert x\rvert^2+\epsilon)$),
block-reduced to $1024$ frequency columns $\times\,8$ time rows per tile, and
normalized with fixed limits $[\,\SI{-22}{},\SI{+35}{}\,]$~dB. Preview and mask
messages are pushed into a process-global, mutex-guarded mailbox with a
backpressure valve that drops the oldest frames; a Holoviz visualizer drains the
mailbox on a $\sim$\SI{30}{\hertz} periodic condition and composites the mask as
a colored overlay (blend $\alpha\!=\!0.7$). This mailbox decouples render cadence
from the detection path so display backpressure never stalls detection. A
separate logging operator reports throughput (MSps/Gbps) and GPU utilization.

\paragraph{GPU-residency summary.}
A single CUDA stream is threaded through the tuple message type
$\langle \text{tensor}, \text{stream}\rangle$ from CHDR $\to$ FFT $\to$ detector,
so cuFFT and every detector kernel run on the same batch stream. Tensors remain
device-resident end to end; the only device-to-host transfers on the hot path are
small scalar mask-population counters.

% ---------------------------------------------------------------------------
\subsection{Coherent-Power Detector: A Hybrid Image/Power Formulation}
\label{sec:detector}

The detector treats the power-dB spectrogram simultaneously as (i) a
\emph{2-D image}, applying local spatial-contrast estimation and binary
morphology, and (ii) a bank of per-frequency \emph{power channels}, applying an
absolute, calibrated noise-floor test. These two views are complementary: the
image view localizes signal boundaries and rejects speckle but hollows out the
interior of wideband signals, whereas the absolute-power view fills wideband
bodies solidly and rescues strong narrowband carriers. The final mask is their
union, cleaned by morphology with a persistence constraint. All stages are custom
CUDA kernels operating in an internal orientation with rows $=$ frequency and
columns $=$ time.

\paragraph{Power-dB formation and front-end equalization.}
From the complex frame the detector computes
$P_{\text{dB}}[f,t]=10\log_{10}\!\big(\Re^2+\Im^2+10^{-12}\big)$. An optional
front-end correction estimates a smoothed per-frequency noise profile
(Gaussian-smoothed row means, $\sigma\!=\!12$ bins) and applies a bounded
per-row gain to flatten the noise floor across frequency---an image
row-normalization---yielding $C_{\text{dB}}$. In the reported real-time profile
the correction gain is clamped to \SI{0}{\decibel} (boost/cap disabled), so
$C_{\text{dB}}\!\approx\!P_{\text{dB}}$ while the machinery remains available for
environments with strong band tilt.

\paragraph{Image view: local spatial support.}
A separable box filter estimates a local background
$B[f,t]$ as the 2-D mean of $C_{\text{dB}}$ over a $\pm$\num{10}-bin time window
and $\pm$\num{8}-bin frequency window. The \emph{support} (local contrast) is
$S[f,t]=C_{\text{dB}}[f,t]-B[f,t]$, normalized and thresholded:
%
\begin{equation}
s[f,t]=\operatorname{clip}\!\Big(\tfrac{S[f,t]-\phi_{\text{floor}}}{\sigma_{\text{span}}},0,1\Big),\qquad
m_{\text{box}}[f,t]=\mathbb{1}\!\left[s\ge\tau\right],
\label{eq:support}
\end{equation}
%
with $\phi_{\text{floor}}=\SI{0.5}{\decibel}$, $\sigma_{\text{span}}=\SI{3}{\decibel}$,
$\tau=0.4$; i.e.\ a pixel is admitted when it exceeds its local neighborhood by
$\ge\SI{1.7}{\decibel}$. This is a locally adaptive (CFAR-like) image detector: it
follows slow gain/noise variation but, because $B$ averages over a wide frequency
window, it high-pass-filters the interior of signals broader than the box.

\paragraph{Power view: absolute per-frequency floor and strong rescue.}
Two absolute-power tests are OR-ed into the mask, each raising the score without
ever lowering it:
\begin{itemize}
  \item \emph{Per-frequency fill.} A calibrated absolute noise floor
  $F[f]$ (one value per frequency bin, loaded from an offline calibration) admits
  any pixel with $C_{\text{dB}}[f,t] > F[f]+\Delta_{\text{fill}}$
  ($\Delta_{\text{fill}}=\SI{2}{\decibel}$). Unlike Eq.~\eqref{eq:support}, this
  references the \emph{global} per-frequency floor rather than a local box mean, so
  the solid interior of a wideband signal is filled instead of hollowed.
  \item \emph{Strong-signal rescue.} A per-frequency floor is re-estimated as the
  time-average of the local background (which dilutes any persistent narrow
  carrier, keeping it a noise estimate). Pixels exceeding this floor by
  $\Delta_{\text{strong}}=\SI{8}{\decibel}$ are flagged, subject to a
  time-persistence guard requiring $\ge\num{3}$ strong time bins within a local
  window (rejecting impulsive spikes). This candidate is OR-ed back in
  \emph{after} the width-based morphology, so a strong but frequency-narrow signal
  survives filters that would otherwise erase it.
\end{itemize}

\paragraph{Image morphology and persistence.}
The unioned binary mask is cleaned with pure image-processing operations: a
$3\times3$ majority filter, a morphological opening ($3\times7$ structuring
element) to remove speckle, a closing ($5\times21$) to bridge gaps, and a
frequency-persistence test that keeps a pixel only if $\ge\num{11}$ of \num{45}
neighboring frequency bins on the same time row are active. The
strong-rescue mask is then OR-ed back in to preserve narrowband detections.

\paragraph{Mask emission.}
The detector emits a full-resolution ($512\times20480$) device-resident binary
mask (\texttt{DetectorMaskMessage}) backed by a recycled device-buffer pool, with
per-frame population statistics attached as metadata for downstream logging and
evaluation. Table~\ref{tab:params} summarizes the active parameters.

\begin{table}[t]
\centering
\caption{Active coherent-power detector parameters (real-time single-channel profile).}
\label{tab:params}
\small
\begin{tabular}{@{}lll@{}}
\toprule
Stage & Parameter & Value \\
\midrule
Front end     & noise-smoothing $\sigma$          & \SI{12}{bins} \\
              & max boost / signal cap            & \SI{0}{\decibel} (disabled) \\
Local background & time / freq box radius         & $\pm10$ / $\pm8$ bins \\
Support test  & floor $\phi_{\text{floor}}$       & \SI{0.5}{\decibel} \\
              & span $\sigma_{\text{span}}$       & \SI{3}{\decibel} \\
              & threshold $\tau$                  & $0.4$ \\
Per-freq fill & offset $\Delta_{\text{fill}}$     & \SI{2}{\decibel} \\
Strong rescue & excess $\Delta_{\text{strong}}$   & \SI{8}{\decibel} \\
              & min.\ time bins                   & $3$ \\
Morphology    & opening / closing SE             & $3{\times}7$ / $5{\times}21$ \\
              & majority filter                   & $3{\times}3$, 1 iter. \\
Persistence   & window / min.\ hits              & $45$ / $11$ bins \\
Sideband guard & ignored per side                & \SI{7}{\mega\hertz} \\
\bottomrule
\end{tabular}
\end{table}

% ---------------------------------------------------------------------------
\subsection{Evaluation Dataset}
\label{sec:dataset}

To evaluate detection across a controlled variety of modern waveforms and channel
conditions, we use over-the-air captures recorded on a USRP~X410 at a sample rate
of \SI{245.76}{\mega\hertz} (complex \texttt{cf32}), storing samples and labels in
the \emph{SigMF} format. Each recording is a remote receive of a locally
transmitted playlist that cycles a library of synthesized waveforms, so that
every emission has a machine-generated ground-truth annotation.

\paragraph{Waveform variety.}
The library spans classical digital modulation, standardized communication
waveforms, analog FM, and synchronization sequences
(Table~\ref{tab:waveforms}). Linear-modulation classes (BPSK/QPSK/16QAM) are
generated across a parametric grid---symbol rates from \SI{3.84}{} to
\SI{61.44}{\mega\hertz}, root-raised-cosine and raised-cosine pulse shaping with
roll-off $\alpha\!\in\!\{0.20,0.35\}$, multiple phase offsets and oversampling
factors---encoded verbatim in a per-annotation \texttt{variation} string
(e.g.\ \texttt{BPSK\_Rs61p44MHz\_occ82p944MHz\_rrc\_a0p35\_pn9}). Across all
classes the library contains \num{287} distinct parameterizations. Standardized
classes include OFDM, 5G downlink, IEEE~802.11ax (Wi-Fi~6), and Bluetooth; analog
classes include narrowband and broadband FM; and Zadoff--Chu (ZC) sequences serve
as fixed-bandwidth synchronization bursts. Dedicated \texttt{METADATA} blocks mark
in-band header regions.

\begin{table}[t]
\centering
\caption{Waveform classes and their occupied-bandwidth ranges in a representative
capture. Emission durations take six discrete lengths
$\{9830,49152,245760,1228800,2457600,4915200\}$ samples
($\approx\SI{40}{\micro\second},\SI{200}{\micro\second},1,5,10,\SI{20}{\milli\second}$
at \SI{245.76}{MSps}); ZC sync bursts are fixed at \num{2955} samples
($\approx\SI{12}{\micro\second}$).}
\label{tab:waveforms}
\small
\begin{tabular}{@{}llr@{}}
\toprule
Class & Type & Occ.\ BW (MHz) \\
\midrule
BPSK / QPSK / 16QAM & Linear PSK/QAM (RRC/RC) & $0.29$--$82.94$ \\
OFDM               & Multicarrier            & $4.50$--$98.28$ \\
5G\_Downlink       & 5G NR DL                & $3.96$--$98.28$ \\
802\_11ax          & Wi-Fi 6                 & $20.0$--$160.0$ \\
Bluetooth          & FHSS / GFSK             & $1.0$--$2.0$ \\
Broadband\_FM      & Analog FM               & $0.13$--$61.47$ \\
Narrowband\_FM     & Analog FM               & $0.01$--$0.03$ \\
ZC                 & Zadoff--Chu sync        & $50.0$ (fixed) \\
\bottomrule
\end{tabular}
\end{table}

\paragraph{Channel conditions.}
The same waveform playlist is captured under a swept transmit attenuation of
$\{0,5,10,\dots,55,60\}$~\si{\decibel} (\num{14} recordings), emulating a range of
received-power / effective-SNR regimes from strong to near-noise-floor. Per-signal
transmit power is held uniform within a capture, so the effective-SNR variation is
captured entirely by this attenuation axis, which lets detection accuracy be
reported as a function of received power and enables apples-to-apples comparison
across detectors and channel conditions.

\paragraph{Metadata schema.}
Ground truth uses SigMF core annotation fields---\texttt{sample\_start},
\texttt{sample\_count}, \texttt{freq\_lower\_edge}, \texttt{freq\_upper\_edge},
and \texttt{label}---augmented with a \texttt{wfgt} (waveform-ground-truth)
extension carrying \texttt{class}, the full parametric \texttt{variation} string,
\texttt{occupied\_bw\_hz}, \texttt{power\_db}, \texttt{length\_samples},
\texttt{block\_center\_hz}, and \texttt{time\_group} (a temporal scheduling index,
enabling per-time-block analysis). Global fields record the datatype, sample rate,
hardware (\texttt{USRP X410}), and channel count. Because every emission carries
its exact time--frequency extent and physical-layer parameters, the same captures
support detection, classification, and future ML studies without re-annotation,
and support portable comparison across systems and environments.

% ---------------------------------------------------------------------------
\subsection{Offline Evaluation Protocol}
\label{sec:eval}

We score detectors offline against the SigMF ground truth so results are exactly
reproducible and detector-agnostic. Each capture is replayed through the identical
signal chain used online---SigMF $\rightarrow$ FFT $\rightarrow$ spectrogram
$\rightarrow$ detector---under a synchronous scheduler with a fresh per-frame
device tensor (eliminating any cross-frame buffer aliasing), so masks stay
time-aligned with the source; a column-profile correlation self-check gates every
run for zero systematic frame offset. The offline FFT geometry is derived from the
SigMF sample rate (\SI{245.76}{MSps} $\rightarrow$ \num{10240}-point FFT,
\SI{24}{\kilo\hertz} bins, \num{512} time rows per frame, $\approx$\num{332}
frames per capture), and is shared by all detectors so comparisons are on a common
grid; each detector's native-resolution mask is nearest-neighbor resampled onto
this common $512\times10240$ grid before scoring.

For each frame we rasterize every annotation into a filled frequency$\times$time
ground-truth rectangle and compare it to the detector mask. We report two
complementary families of metrics: (i) \emph{pixel-level} precision, recall, F1,
IoU, and false-positive area, characterizing how clean and well-localized
detections are; and (ii) \emph{region-level} coverage (fraction of each
ground-truth box covered by the mask) and detection rate at a coverage threshold,
characterizing whether each emission was found. Because ground truth is a filled
box while energy-based detectors fire sparsely on signal support, we emphasize
region detection rate and mean coverage for ``did it find the signal'' and pixel
precision for ``are detections clean.'' Every metric is broken down along
registered dimensions---signal class, occupied bandwidth, pulse length, transmit
attenuation (the effective-SNR axis), and time group---so accuracy can be reported
per waveform family and per channel condition. Maskless tail frames dropped to pipeline drain
are recorded but excluded from aggregates so they do not bias the comparison.
